\section{Compact equations from an existing denominator}
\label{sec:compact}

The second contribution is a family of explicit polynomials in ten
strictly positive integer unknowns.  It improves the latest incoming
\(78/630\) and \(79/468\) circuits, where the two numbers mean live binary
arithmetic operations and exact total degree.  The new arithmetic permits
an existing expression to serve as the carrier denominator; two further
savings reuse products already needed elsewhere.  The complete results
are as follows.

\begin{center}
\begin{tabular}{@{}llrr@{}}
\toprule
Denominator & Transport weight & Operations & Exact degree\\
\midrule
\(XY\) & \(V^2\) & \textbf{75} & \textbf{626}\\
\(XY\) & \(bV\) & \textbf{76} & \textbf{464}\\
\(q^4\) & \(V^2\) & 76 & 622\\
\(q^4\) & \(bV\) & \textbf{77} & \textbf{460}\\
\bottomrule
\end{tabular}
\end{center}

The \(76/622\) point has the simpler squared-weight proof, although
\(76/464\) dominates it in the two numerical metrics.  The other three
points trade arithmetic cost against degree.  All four use the same
complete indexed BFS943 coefficient tuple, specified in
Section~\ref{sec:compact-constants}; its four ports are constants, not
additional unknowns.  The arithmetic changes are uniform in those ports
within their stated box, but the consistency interpretation uses the
calibration of that \emph{whole} tuple.\cite{certificate,indexed}

\begin{theorem}[Compact consistency equations]
\label{thm:compact-main}
Each of the four fixed polynomials defined below has ten strictly positive
integer unknowns.  Composing the ordinary arithmetic proof in this
section with the retained indexed native/refutation calibration gives
\[
 \ZFC\ \vdash\
 \bigl[\nexists\,\mathbf x\in\Z_{>0}^{10}:P(\mathbf x)=0\bigr]
 \ \longleftrightarrow\ \Con(\ZFC).
\]
Their operation counts and exact total degrees are those in the table.
\end{theorem}

This is an ordinary theorem formalizable in ZFC.  It is not a new
end-to-end Lean or Rocq derivation, an emitted object-language ZFC proof,
a proof of consistency, or a claim of global optimality.  Pell doubling,
the product-square finalizer, and the removal of the auxiliary shift have
explicit predecessors.  The new denominator criterion and its completed
composition are distinguished from those inherited steps below.
\cite{certificate,unshift}

\subsection{The complete polynomial}
\label{subsec:compact-definition}

The independent unknowns are exactly
\[
 J,F,\alpha,i,T,w,s,n,\eta,\zeta\in\Z_{>0}.
\]
The four fixed positive integers are \(\beta,K_0,M_C,M_F\), with
\begin{equation}
 \beta\ge31\text{ odd},\qquad K_0>0,\qquad
 0<M_C<\beta<M_F<2\beta.
 \label{eq:compact-portbox}
\end{equation}
Define the following polynomial abbreviations:
\begin{align}
 h&=\beta J,& q&=h+1,& C&=q-F-\alpha,\nonumber\\
 X&=wq,&Y&=q^3s,&k&=\eta+\zeta,\nonumber\\
 c&=kY+\eta,& A&=Y(X+1)+1,&
 \Delta&=A^2-1,\nonumber\\
 H&=4A-5,& \ell&=(XY)(kY),&
 E&=1+\ell(\ell+k).                         \label{eq:compact-base}
\end{align}
For the mask, transport, and strong norm, put
\begin{align}
 R&=[h(q-F)+\alpha](q^2-1)+(M_C+qM_F)J,\nonumber\\
 e&=k-R-1, & B&=(K_0+w)C-F,\nonumber\\
 U&=ic^2,& a&=\Delta U^2E,& F_0&=1+a,\nonumber\\
 G&=EF_0,& \mathcal F&=1+4aF_0,&
 Q&=\Delta(\mathcal F-1),\nonumber\\
 b&=cT-R,& V&=\mathcal F b-c,&
 K_+&=X(3c-2X)+2.                            \label{eq:compact-strong}
\end{align}
Here \(\mathcal F=(1+2a)^2\) is an identity; its square root has not
been introduced as an unknown.  Choose one of
\[
 D=XY\quad\hbox{or}\quad D=q^4,
 \qquad
 L=V^2\quad\hbox{or}\quad L=bV.
\]
Finally set
\begin{align}
 M&=hDH,\nonumber\\
 Z&=hD\bigl[H(n-LB)+K_+\bigr]-H(DB+he),\nonumber\\
 R_G&=M^2G.                                             \label{eq:compact-carrier}
\end{align}
The requested single equation is
\begin{equation}
 \boxed{\ P=Q(Z^2-R_GV^2)+(R_G-Z^2)=0.\ }
 \label{eq:compact-final}
\end{equation}
Every named quantity in these displays is a polynomial abbreviation.
The delivered JSON graphs contain only binary addition, subtraction, and
multiplication.  Squaring and multiplication by fixed scalars are counted
as multiplications.  No division or variable exponent occurs in the
polynomial.

The displayed form also makes the two shared products visible:
\(A-1=XY+Y\), while \(kY\), already used in \(c\), is reused in
\(\ell=(XY)(kY)\).  The proof occasionally writes \(V_0=XY^2\), but
the evaluator need not compute that expression separately.

\subsection{The admissible-denominator criterion}
\label{subsec:compact-denominator}

The same proof works with a more general integer-polynomial denominator
\(D\), provided that, at every positive input,
\begin{equation}
 D>0,\qquad q^4\mid D,\qquad D\mid4XY^2.
 \label{eq:compact-admissible}
\end{equation}
The first divisibility will recover a bounded index error.  The second
will recover the complete quotient only \emph{after} that error is known
to vanish.  Pairwise coprimality of \(D\) with all the other denominators
is unnecessary.  Both advertised choices satisfy the criterion:
\begin{align*}
 XY&=q^4(ws),&4XY^2&=(XY)(4Y),\\
 4XY^2&=q^4(4ws^2q^3).
\end{align*}
In particular, \(D\ge q^4>2\).

Both choices depend only on the outer coordinates \(J,w,s\), so the
canonical completeness construction below has no circular choice of a
padding coordinate.  In fact the pointwise criterion already excludes
such a dependence for any polynomial \(D\): with \(J,w,s\) fixed, it
implies \(0<D\le4XY^2\) while any of the other coordinates varies through
all positive integers.  A bounded polynomial in such a coordinate is
constant.  Repeating this argument removes every non-outer coordinate.
Within this criterion, total degree four is minimal for a polynomial
denominator: fix \(w,s\) and let \(J\) grow, using \(D\ge(\beta J+1)^4\).
The choice \(q^4\) attains that denominator-degree bound; this is not a
global bound on consistency equations.

The rational quantities used only in the proof are
\[
 z=\frac ZM,\qquad \kappa=\frac{K_+}{H},\qquad
 \rho=\frac eD,\qquad t=\frac Bh,\qquad W=1+hL.
\]
The exact identities are
\begin{align}
 z&=n+\kappa-Wt-\rho,                         \label{eq:compact-rational}\\
 P&=(Q-1)Z^2-M^2G(QV^2-1).                   \label{eq:compact-norm}
\end{align}
No divisibility premise is used in either identity.  In particular,
\(\rho\) and \(t\) are not initially assumed integral.

\subsection{Pell facts and the index lemmas}
\label{subsec:compact-pell}

For an integer \(A\ge2\), define
\[
 \chi_A(r)+\psi_A(r)\sqrt{A^2-1}
       =(A+\sqrt{A^2-1})^r\qquad(r\ge0).
\]
Both sequences satisfy \(u_{r+2}=2Au_{r+1}-u_r\), with
\(\chi_A(0)=1,\chi_A(1)=A\) and
\(\psi_A(0)=0,\psi_A(1)=1\).
Positive integer solutions of \(u^2-(A^2-1)v^2=1\) are exactly these
pairs.  One elementary descent proves this classification: for \(v>1\),
multiplication by \(A-\sqrt{A^2-1}\) gives a positive integer ordinate
\(Av-u\) strictly smaller than \(v\); at \(v=1\), \(u=A\).
The addition formulas also give strong divisibility,
\begin{equation}
 \gcd\bigl(\psi_A(r),\psi_A(t)\bigr)
       =\psi_A(\gcd(r,t)).                         \label{eq:compact-gcd}
\end{equation}
Indeed, \(\chi_A(r)\) is a unit modulo \(\psi_A(r)\), and the addition
law reduces the gcd of two ordinates exactly as the Euclidean algorithm
reduces their indices.  These elementary facts are the Pell framework
of the retained normalized kernel.\cite{native61}

We also need the first norm in \eqref{eq:compact-base}.  If \(v\ge1\)
and positive integers \(\tau,k\) satisfy
\(\tau^2-v(v+1)k^2=1\), then
\begin{equation}
 \tau=\chi_{2v+1}(r),\qquad k=2\psi_{2v+1}(r)
 \quad\hbox{for some }r\ge1.                       \label{eq:compact-first-class}
\end{equation}
For completeness, there is no solution with \(k=1\), because
\(v(v+1)+1\) lies strictly between consecutive squares.  The pair
\((2v+1,2)\) is a solution.  For \(k>2\), multiplication by its inverse
unit gives the positive smaller ordinate \((2v+1)k-2\tau\).  Descent
therefore classifies every positive solution as a power of this pair,
which is precisely \eqref{eq:compact-first-class}.

\begin{lemma}[Odd-index quotient identities]
\label{lem:compact-odd}
There are integer polynomials \(Q_r\) such that
\begin{align*}
 \chi_t(2r+1)&=tQ_r(t^2),&
 Q_r(0)&=(-1)^r(2r+1),\\
 Q_r(1-A^2)&=(-1)^r\psi_A(2r+1),&
 \psi_t(2r+1)&\equiv(-1)^r\pmod{t^2}.
\end{align*}
\end{lemma}
\begin{proof}
Use \(Q_0(x)=1\), \(Q_1(x)=4x-3\), and
\(Q_{r+1}(x)=(4x-2)Q_r(x)-Q_{r-1}(x)\).
The odd trace subsequence gives the first identity.  Evaluation at zero
gives the second.  At \(x=1-A^2\), the recurrence coefficient is the
negative of the one for the odd \(\psi_A\) subsequence, whose initial
values are \(1,4A^2-1\); this gives the third identity.  The recurrence
for \(\psi_t(2r+1)\), with initial values \(1,4t^2-1\), gives the last
congruence modulo \(t^2\).
\end{proof}

\begin{lemma}[Normalized rank]
\label{lem:compact-rank}
Suppose \(c=\psi_A(p)>2\), \(f,I>0\),
\(f^2-(A^2-1)I^2=1\), and \(c^2\mid I\).
Then \(f=\chi_A(m)\), \(I=\psi_A(m)\) for an index satisfying
\(pc\mid m\), and hence \(c\mid m\) and \(m>2p\).
\end{lemma}
\begin{proof}
Classification gives \(m\), and \eqref{eq:compact-gcd} gives \(p\mid m\).
Write \(m=pv\) and \(D_0=\chi_A(p)\).  Composition gives
\(I=c\psi_{D_0}(v)\).  Since \(D_0^2\equiv1\pmod c\), recurrence
gives \(\psi_{D_0}(v)\equiv vD_0^{v-1}\pmod c\).
The number \(D_0\) is invertible modulo \(c\), and \(c^2\mid I\), so
\(c\mid v\).  Thus \(pc\mid m\) and \(m\ge pc>2p\).
\end{proof}

\begin{lemma}[Signed index recovery]
\label{lem:compact-index}
Let \(c=\psi_A(p)\), \(f=\chi_A(m)\), \(c\mid m\), \(m>2p\),
\(S\ge2\), \(c\mid S\), and \(S^2=(A^2-1)(f^2-1)\).
If \(0<p,R<c/2\) and integers \(V,y\) satisfy
\[
 S^2V^2-(S^2-1)y^2=1,\qquad
 V\equiv-c\pmod f,\qquad V\equiv-R\pmod c,
\]
then \(p=R\).  The signs of \(V,y\) are unrestricted.
\end{lemma}
\begin{proof}
Classification at parameter \(S\) gives
\(S|V|=\chi_S(\nu)\), \(|y|=\psi_S(\nu)\).
The index \(\nu\) is odd: an even trace is \(\pm1\pmod S\).
Lemma~\ref{lem:compact-odd}, using
\(S^2\equiv1-A^2\pmod f\), gives
\(\psi_A(\nu)\equiv\pm c\pmod f\).  Squaring and using
\(\chi_A(2j)=1+2(A^2-1)\psi_A(j)^2\) therefore gives
\[
 \chi_A(2\nu)\equiv\chi_A(2p)\pmod{\chi_A(m)}.
\]
Here is the abscissa step used to read this congruence.  For
\(0<k<m\), a congruence
\(\chi_A(t)\equiv\chi_A(k)\pmod{\chi_A(m)}\) implies
\(t\equiv\pm k\pmod{2m}\).  The addition laws reduce and reflect
\(t\) to a signed representative \(\pm\chi_A(u)\), \(0\le u\le m\).
The value \(u=m\) has residue zero.  The negative sign is impossible
because \(0<\chi_A(u)+\chi_A(k)\le2\chi_A(m-1)<\chi_A(m)\).
The positive sign forces equality and hence \(u=k\).
Apply this step with \(k=2p<m\): it gives
\(\nu\equiv\pm p\pmod m\), hence modulo \(c\).

On the other hand, \(c\mid S\) and \(Q_r(0)=(-1)^r(2r+1)\), together
with \(V\equiv-R\pmod c\), give \(\nu\equiv\pm R\pmod c\).
Thus \(p\equiv\pm R\pmod c\).  The strict half-interval bounds on
\(p,R\) force \(p=R\).
\end{proof}

\subsection{Root extraction before any transport sign}
\label{subsec:compact-roots}

At every positive input,
\[
 q\ge32,\quad X\ge q,\quad Y\ge q^3,\quad k\ge2,\quad c\ge3,
 \quad A>4,\quad\Delta\ge3,\quad U\ge c^2,\quad E>1.
\]
Moreover \(\mathcal F>2c\), \(Q>1\), and \(M>0\).
For every integral \(b\), \(V=\mathcal Fb-c\) is nonzero and
\(|V|>1\): if \(b=0\), \(V=-c\); otherwise
\(|V|\ge\mathcal F-c>1\).
The identities
\begin{equation}
 \mathcal F=(1+2a)^2,\qquad
 Q=4\Delta^2U^2G                                      \label{eq:compact-double}
\end{equation}
imply \(\gcd(G,Q-1)=\gcd(G,QV^2-1)=1\).
At a zero of \(P\), equation~\eqref{eq:compact-norm} yields
\[
 (Q-1)z^2=G(QV^2-1).
\]
For every prime dividing \(G\), its valuations in \(Q-1\) and
\(QV^2-1\) vanish; its valuation in \(G\) must therefore be even,
because the valuation of a rational square is even.  Thus \(G=g^2\)
for an integer \(g>0\).  Since \(F_0\equiv1\pmod E\), the coprime
positive factors \(E,F_0\) are separately squares:
\[
 E=\tau^2,\qquad F_0=f_0^2,\qquad g=\tau f_0.
\]
No coprimality condition on \(M\) was required.
Put \(y=|z|/g\).  Then
\begin{equation}
 QV^2-(Q-1)y^2=1.                                  \label{eq:compact-adjacent}
\end{equation}

\begin{lemma}[Integral ordinate in the adjacent norm]
\label{lem:compact-adjacent}
For an integer \(Q>1\), a nonzero integer \(V\), and a positive rational
\(y\), equation~\eqref{eq:compact-adjacent} forces \(y\) integral.
\end{lemma}
\begin{proof}
Write \(y=r/m\) in lowest terms.  Clearing denominators gives
\(m^2\mid Q-1\).  For \(v=|V|>1\), the norm implies
\(v<y<Qv/(Q-1)\).  The transformation
\[
 v'=(2Q-1)v-2(Q-1)y,\qquad
 y'=(2Q-1)y-2Qv
\]
preserves the norm.  Its first coordinate is integral and satisfies
\(-v<v'<v\); it cannot vanish, since the norm is one.
The reduced denominator of the second coordinate remains \(m\), because
\(2Q-1\equiv1\pmod m\).  Reflecting signs and descending on the positive
integer first coordinate eventually reaches \(v=1\), where the norm
gives \(y=1\).  Hence the unchanged denominator is \(m=1\).
\end{proof}

Consequently \(z=\varepsilon gy\) is an integer, with
\(\varepsilon\in\{-1,1\}\).  For proof purposes set
\begin{equation}
 f=1+2a,\qquad I=2Ug.
 \label{eq:compact-proofpair}
\end{equation}
Then
\begin{equation}
 f^2=\mathcal F,\qquad f^2-\Delta I^2=1,\qquad
 Q=(\Delta I)^2,\qquad c^2\mid I.                 \label{eq:compact-strongnorm}
\end{equation}

We next give the sign-free main-index folding argument explicitly.
Classify \(f=\chi_A(m)\), \(I=\psi_A(m)\).
At \(S=\Delta I\), classify \eqref{eq:compact-adjacent} as
\(S|V|=\chi_S(j)\), \(y=\psi_S(j)\).  Again \(j\) is odd.
Since \(S^2\equiv1-A^2\pmod f\),
Lemma~\ref{lem:compact-odd} and \(V\equiv-c\pmod f\) give
\(\psi_A(j)\equiv\pm c\pmod f\).  The addition laws give
\[
 \psi_A(j+2m)\equiv-\psi_A(j),\qquad
 \psi_A(2m-j)\equiv\psi_A(j)\pmod f.
\]
Reduction modulo \(2m\), followed by reflection about \(m\), supplies
\(p\in[0,m]\) of the parity of \(j\), with
\(\psi_A(j)\equiv\pm\psi_A(p)\pmod f\).
Now \(0<c\le I\), \(0\le\psi_A(p)\le I\), and \(2I<f\) because
\(A\ge3\).  The congruence therefore forces
\begin{equation}
 c=\psi_A(p),\qquad p\ge1\text{ odd}.              \label{eq:compact-folded}
\end{equation}
Neither the sign of \(V\) nor a parity condition on \(m\) was imposed.

Write \(V_0=XY^2\) for proof notation.  The first norm and
\eqref{eq:compact-first-class} give
\[
 k=2\psi_{2V_0+1}(n_0),\qquad n_0\ge1.
\]
If \(n_0=1\), then \(k=2\), \(\eta=\zeta=1\), and \(c=2Y+1\)
lies strictly between \(\psi_A(1)=1\) and \(\psi_A(2)=2A\), a
contradiction.  Thus
\begin{equation}
 n_0\ge2,\quad k\ge8V_0+4,\quad
 c>kY\ge k>8q^7,\quad c>X,\quad K_+>0.           \label{eq:compact-earlybound}
\end{equation}
No sign of \(B,C,R,b,V\), or \(\rho\) was used.

\subsection{Signed transport with no shift in \texorpdfstring{$T$}{T}}
\label{subsec:compact-sign}

The old coordinate was \(T+\alpha\).  Its formal inverse after deletion
is \(T_{\rm old}=T-\alpha\), which can be nonpositive.  Soundness therefore
needs a direct argument; a formal substitution alone would be inadequate.
\cite{unshift}

First \(b\ne0\).  Otherwise \(V=-c\), and
\eqref{eq:compact-adjacent} would give
\[
 y^2-c^2=\frac{c^2-1}{Q-1}\in(0,1),
\]
because \(Q\ge9c^4>c^2\).  The left side is integral.
Put \(D_0=\Delta U^2\ge243\).  Since \(a=D_0E\) and
\(g^2=E(1+a)\),
\[
 \frac{a^2}{g^2}=\frac{D_0a}{1+a}>\frac{D_0}{2}>9.
\]
Hence \(a>3g\), while \(a\ge c\).  For nonzero integral \(b\),
\begin{equation}
 \operatorname{sgn}V=\operatorname{sgn}b,\qquad
 |V|\ge\mathcal F-c>a>3g,\qquad |V|>|b|,\qquad
 y<\tfrac32|V|.                                    \label{eq:compact-absolute}
\end{equation}
For \(L=V^2\), these immediately give \(gy<L/2\).

The choice \(L=bV\) requires a stronger estimate.  Suppose the odd
main index in \eqref{eq:compact-folded} were \(p=3\), so
\(c=4A^2-1\).  If \(n_0\ge3\), then
\[
 k>32V_0^2,\qquad
 c<64X^2Y^2<32X^2Y^5<kY<c,
\]
using \(A<4XY\) and \(Y\ge q^3\), a contradiction.
Thus \(n_0=2\) and \(k=8V_0+4\).
Writing \(u=X+2-2Y\), the exact remainder identity is
\begin{equation}
 \eta=4Y^2(Xu+1)+8Y(X+1)+3-4Y.                    \label{eq:compact-p3}
\end{equation}
If \(u\le-1\), this is negative because
\(Y(X-1)>2(X+1)\).  If \(u\ge2\), subtraction of \(k\) leaves at
least \(4Y^2+8Y(X+1)-4Y-1>0\).  Since \(0<\eta<k\), one must have
\(u\in\{0,1\}\), equivalently \(2Y-X\in\{1,2\}\).  But
\(q\mid X,Y\) and \(q\ge32\).  This contradiction proves
\begin{equation}
 p\ge5,\qquad c>c_3:=4\Delta+3.                   \label{eq:compact-p5}
\end{equation}

At \(S=\Delta I\), the adjacent index \(j\) is odd.  Index one gives
\(|V|=1\), already excluded.  Index three gives \(|V|=4Q-3\).
The two possible signs would imply, respectively,
\[
 \mathcal F(b-4\Delta)=c-c_3,\qquad
 \mathcal F(b+4\Delta)=c+c_3.
\]
Both right sides lie strictly between zero and \(\mathcal F\), since
\(0<c-c_3<c+c_3<2c<\mathcal F\).  Thus neither can be a multiple of
\(\mathcal F\).  Consequently \(j\ge5\) and
\[
 |V|\ge16Q^2-20Q+5>8Q^2.
\]
Since \(Q>\mathcal F>2c\), the triangle inequality gives
\[
 |b|>\frac{8Q^2-c}{\mathcal F}
       >8\mathcal F-\tfrac12>4\mathcal F>4g.
\]
Together with \eqref{eq:compact-absolute}, this proves
\begin{equation}
 L=bV=|b||V|>0,\qquad gy<\tfrac38L<L/2.            \label{eq:compact-bvbound}
\end{equation}

For either weight, the new unshifted identity is
\[
 e=b-cT+k-1<b,\qquad \rho<|b|/D,
\]
using \(c>k\) and \(T\ge1\).  If \(B\le-1\), then
\(Wt\le-L-1/h\).  Equation~\eqref{eq:compact-rational} gives
\[
 n=\varepsilon gy-\kappa+Wt+\rho
     <L/2-L+|b|/D<0.
\]
The last inequality uses \(D>2\) and either \(L=V^2\) with
\(|V|>|b|\), or \(L=|b||V|\) with \(|V|>1\).
This contradicts the positive independent input \(n\).  Hence
\begin{equation}
 B\ge0,\qquad C>0.                               \label{eq:compact-transportpos}
\end{equation}

\subsection{Rank recovery and the staged denominator argument}
\label{subsec:compact-staged}

The outer bounds can now be proved without source semantics.  Let
\(G_0=h(q-F)+\alpha\).  As \(q-F=C+\alpha\),
\[
 G_0=q\alpha+(q-1)C,\qquad q\le G_0\le q^2-q,
\]
while \eqref{eq:compact-portbox} gives
\(0<(M_C+qM_F)J<h(1+2q)\).  Therefore
\begin{equation}
 3q<R<q^4-2.                                      \label{eq:compact-outerbound}
\end{equation}
From \eqref{eq:compact-earlybound}, \(c>2R\).  Also
\(c=\psi_A(p)>2p\): \(p\ge3\), and the Pell recurrence gives more
than geometric growth at \(A\ge3\).  Lemma~\ref{lem:compact-rank}
supplies \(m\) with \(c\mid m\), \(m>2p\).
Since \(c^2\mid I\), the strong norm gives \(f^2\equiv1\pmod c\), so
\[
 V\equiv-c\pmod f,\qquad V\equiv-R\pmod c.
\]
Lemma~\ref{lem:compact-index} applies with \(S=\Delta I\), giving
\begin{equation}
 p=R.                                             \label{eq:compact-rankexact}
\end{equation}
In addition, \(T\ge1\) and \(c>2R\) restore \(b\ge c-R>0\) and
\(V>0\).  No later step needs the stronger, generally false inequality
\(b>c\).  The retained mask parity argument yields even \(q\) and odd
\(J\), but the following denominator cancellation does not require that
parity conclusion.\cite{native61,certificate}

Because \(z\) is integral, multiplying
\eqref{eq:compact-rational} by \(hH\) gives
\begin{equation}
 D\mid hHe.                                      \label{eq:compact-partialdiv}
\end{equation}
But \(H=4Y(X+1)-1\equiv-1\pmod q\) and \(h=q-1\), so
\(\gcd(q,hH)=1\).  Since \(q^4\mid D\),
\begin{equation}
 q^4\mid e.                                      \label{eq:compact-smallmodulus}
\end{equation}
The first norm has parameter \(P_0=2XY^2+1\).
The base-one recurrence gives
\(\psi_{P_0}(n_0)\equiv n_0\pmod{P_0-1}\); hence
\begin{equation}
 k\equiv2n_0\pmod{4XY^2}.                        \label{eq:compact-firstcong}
\end{equation}
Equations~\eqref{eq:compact-smallmodulus}--\eqref{eq:compact-firstcong}
imply \(q^4\mid2n_0-R-1\).

Now \(P_0>A\), since
\(P_0-A=Y(2XY-X-1)>0\).  If \(n_0\ge R\), monotonicity in the
positive index and base gives
\[
 k/2=\psi_{P_0}(n_0)\ge\psi_A(R)=c>k,
\]
impossible.  Thus \(1\le n_0<R\), and
\[
 |2n_0-R-1|<R<q^4.
\]
The bounded multiple of \(q^4\) must vanish:
\begin{equation}
 2n_0=R+1.                                       \label{eq:compact-firstindex}
\end{equation}
Only now return to the \emph{full} recurrence modulus in
\eqref{eq:compact-firstcong}.  It gives
\[
 4XY^2\mid e,\qquad D\mid e,\qquad \rho\in\Z.
\]
The middle implication uses the second half of
\eqref{eq:compact-admissible}.  No assumption that \(D\) itself is
coprime to \(hH\) has entered the argument.

\subsection{The retained projection and native calibration}
\label{subsec:compact-transfer}

Write \(R=2r+1\); equation~\eqref{eq:compact-outerbound} gives \(r\ge48\).
The elementary recurrence bounds
\[
 \psi_A(j)\ge(2A-1)^{j-1},\qquad
 \psi_{P_0}(j)\le(2P_0)^{j-1}\qquad(j\ge1)
\]
give a convenient self-contained height estimate.  Indeed,
\[
 (2A-1)^2-2XP_0
 =4Y^2(2X+1)+4Y(X+1)+1-2X>0,
\]
so
\[
 \frac ck=
 \frac{\psi_A(2r+1)}{2\psi_{P_0}(r+1)}
 \ge\frac12\left(\frac{(2A-1)^2}{2P_0}\right)^r
 >\frac{X^r}{2}.
\]
Since \(Y<c/k<Y+1\) and \(X^r>6\), this proves
\begin{equation}
 Y>X^r/3,\qquad X^{r+1}<3Y(X+1)<H.                \label{eq:compact-ratiobound}
\end{equation}

After \(\rho\) is integral, the carrier gives \(H\mid hK_+\).
The needed projection can also be proved directly.  Modulo the odd
integer \(H=4A-5\), put \(u_0=0\) and
\(u_j=3\cdot2^{j-1}\psi_A(j)\) for \(j\ge1\).
The Pell recurrence and \(4A\equiv5\pmod H\) give
\[
 u_{j+2}\equiv5u_{j+1}-4u_j\pmod H,\qquad u_0=0,\quad u_1=3.
\]
The sequence \(4^j-1\) has the same recurrence and initial values.
Thus, writing \(x_*=2^R\),
\[
 3x_*c\equiv2(x_*^2-1)\pmod H,\qquad
 x_*K_+\equiv-2(X-x_*)(x_*X+1)\pmod H.
\]
Multiplying the known divisibility by \(x_*\) and cancelling only the
unit two gives
\begin{equation}
 H\mid h(X-2^R)(2^RX+1).                         \label{eq:compact-projection}
\end{equation}
There is no division by three in this argument.  For the height of the
\emph{entire product}, use \(h<X\), \(X+x_*<Xx_*\), and
\(x_*X+1<2x_*X\):
\[
 \bigl|h(X-x_*)(x_*X+1)\bigr|
 <2x_*^2X^3=8\cdot16^rX^3<X^{r+1}<H.
\]
The penultimate inequality follows from \(X\ge32\) and \(r\ge48\)
(already \(r>13\) is enough).  The product must vanish, and hence
\begin{equation}
 X=2^R.                                          \label{eq:compact-power}
\end{equation}
This is not a cancellation of possibly nonunit individual factors modulo
a composite integer.  Since \(x_*\) is a power of two and \(H\) is odd,
the same congruence now gives \(H\mid K_+\), so \(\kappa\in\Z\).
The carrier then makes \(Wt\) integral; since
\(W\equiv1\pmod h\), \(h\mid B\), and \(t\in\Z_{\ge0}\).

At this point the remaining inherited interface is explicit:
the fixed-minus parity, exact half-binomial rounding, and indexed native
calibration are retained from the supplied ordinary sources.  They are
not inferred from the finite circuit checks.\cite{certificate,native61,indexed}
All hypotheses of that retained fixed-minus parity, exact half-binomial
rounding, and native reconstruction are now available, including the
literal \(q^3\mid Y\).  Reconstruction uses the fresh source base
\(Y(X+1)+2\); it does not identify that base with the present \(A\).
The whole indexed source calibration sends these data to a finite
syntactic ZFC refutation.  This establishes soundness at the explicitly
retained ordinary calibration boundary.

\subsection{Positive completeness}
\label{subsec:compact-completeness}

Conversely, a finite refutation supplies the inherited canonical outer
data, with \(X=2^R\), \(R>3q\), \(R\equiv3\pmod4\),
\(q^3\mid Y\), and an integer \(t=B/h\ge0\).
The constructive scale-one ratio theorem supplies
\begin{align*}
 c&=\psi_A(R),&
 k&=2\psi_{P_0}((R+1)/2),&
 \tau&=\chi_{P_0}((R+1)/2),\\
 \eta&=c-kY>0,&\zeta&=k-\eta>0.
\end{align*}
Here \(P_0=2XY^2+1\).  The first-index congruence gives
\(4XY^2\mid e\), so \eqref{eq:compact-admissible} gives
\(\rho=e/D\in\Z_{\ge0}\).  Projection gives positive integral
\(\kappa\).

The Pell matrix
\[
 \begin{pmatrix}A&\Delta\\1&A\end{pmatrix}
\]
has determinant one, and hence finite order modulo \(c^2\tau\).
A positive multiple of this period supplies a seed pair \((f_0,I_0)\)
with \(c^2\tau\mid I_0\).  Put \(i=I_0/(c^2\tau)>0\), so
\(U=I_0/\tau\).  Then
\[
 a=\Delta I_0^2=f_0^2-1,\qquad
 F_0=f_0^2,\qquad g=\tau f_0.
\]
Doubling this single seed pair gives
\[
 f=2f_0^2-1=1+2a,\qquad
 I=2f_0I_0=2Ug,
\]
exactly the proof pair \eqref{eq:compact-proofpair}.  There is no
simultaneous-solvability assertion about unrelated Pell equations.

The auxiliary witnesses can also be written directly.  With
\(S=\Delta I\) and \(R=2v+1\), where \(v\) is odd, set
\[
 V_*=\frac{\chi_S(R)}{S}=Q_v(S^2),\qquad
 y_* =\psi_S(R).
\]
They are positive integers satisfying the adjacent norm.  Since
\(S^2\equiv1-A^2\pmod{f^2}\) and \(c\mid S\),
Lemma~\ref{lem:compact-odd} gives
\[
 V_*\equiv-c\pmod{f^2},\qquad V_*\equiv-R\pmod c.
\]
The strong norm gives \(f^2\equiv1\pmod c\) and
\(\gcd(c,f^2)=1\).  Therefore
\begin{equation}
 T=\frac{V_*+c+Rf^2}{cf^2}\in\Z_{>0},\qquad
 f^2(cT-R)-c=V_* .                              \label{eq:compact-Tcomplete}
\end{equation}
In particular, no subtraction of \(\alpha\) is needed.

For either prescribed positive weight \(L\), choose
\[
 n=gy_*+(1+hL)t+\rho-\kappa.
\]
All summands are integers.  The inequalities
\[
 f_0>(A-1)I_0>Xc>\kappa
\]
show \(g>\kappa\); the last inequality follows already from
\(H>3\) and \(K_+<3Xc\).  Since \(y_*\ge1\), \(t,\rho\ge0\),
this gives \(n>0\).  Its carrier is exactly \(z=gy_*\), so
\eqref{eq:compact-final} vanishes.  Every independent coordinate is
strictly positive.  This proves completeness by fresh canonical
witnesses and completes the proof of Theorem~\ref{thm:compact-main}.

\subsection{Exact degrees, arithmetic costs, and finite evidence}
\label{subsec:compact-metrics}

Removing the private \(T+\alpha\) addition saves one operation.
Sharing \(XY\) in the main base and first norm saves a multiplication
when \(D=XY\).  Reusing the already computed \(kY\) in
\(\ell=(XY)(kY)\) saves another multiplication.  The \(q^4\) option
then adds the single square \((q^2)^2\).  Counts were obtained from
the complete live graphs, including fixed scalar multiplications:
\begin{center}
\begin{tabular}{@{}llrrrr@{}}
\toprule
\(D\)&\(L\)&Products&Sums&Differences&Total\\
\midrule
\(XY\)&\(V^2\)&44&17&14&75\\
\(XY\)&\(bV\)&45&17&14&76\\
\(q^4\)&\(V^2\)&45&17&14&76\\
\(q^4\)&\(bV\)&46&17&14&77\\
\bottomrule
\end{tabular}
\end{center}

Write \(q_0=\beta J\), \(k_0=\eta+\zeta\).
The full highest homogeneous components, in table order, are
\begin{align}
 &2^{14}q_0^{378}w^{48}s^{110}i^{20}k_0^{64}T^4
       (q_0-F-\alpha)^2,\nonumber\\
 &-2^{10}q_0^{284}w^{36}s^{82}i^{14}k_0^{46}T^2,\nonumber\\
 &2^{14}q_0^{378}w^{46}s^{108}i^{20}k_0^{64}T^4
       (q_0-F-\alpha)^2,\nonumber\\
 &-2^{10}q_0^{284}w^{34}s^{80}i^{14}k_0^{46}T^2.
 \label{eq:compact-leading}
\end{align}
Their degrees are \(626,464,622,460\).  The selected coefficients
\(2^{14}\beta^{378}\) and \(-2^{10}\beta^{284}\) are nonzero at the
actual fixed \(\beta\).  These are complete multivariate leading
components, not numerical rays.  The squared-weight outputs have 390
leading monomials; the \(bV\) outputs have 47.

The supplied standard-library checker proves 32 or 33 exact
free-polynomial cut identities covering every actual gate, checks full
topology and liveness, collects every leading component, and validates
the four literal coefficient payloads and bit lengths.  A separate
reconstruction script regenerates the integers from the finite code in
Section~\ref{sec:compact-constants}, byte for byte.  Normal and optimized
Python runs agree.  A separately tasked AI reviewer checked the ordinary
proof and independently verified the complete leading forms and graph
liveness; its 480 signed modular scalar comparisons are additional
regression evidence, not substitutes for the exact identities or the
unbounded proof.  The independent review led to the explicit folding
proof printed above.

The emitted JSON graphs, coefficient bytes, builders, checkers,
receipts, source ledger, and copied ordinary proof sources are included
in the archive's \texttt{compact/} directory.  They specify actual fixed
polynomials.  The native/refutation theorem remains the expressly
retained calibration; the finite programs do not establish consistency
or replace that theorem.
